% !PRINTED 214
\weberpage{214}
%<<CONT
Gruppe $\mathfrak{G}$ vom Grade $\mu$ gewiß in $\mathfrak{A}$ enthalten ist. Es ergibt sich hieraus

\weq{16}{\mu = \nu\varphi(n)}

und zugleich die Identität von $\mathfrak{G}$ mit $\mathfrak{A}$.

Daraus folgt aber auch, daß $F(t)$ dem Produkt der sämtlichen Faktoren $F(t,\varrho^m)$ gleich ist, also

\weq{17}{F(t) = \overset{m}{\varPi} F(t, \varrho^m),}

und da $F(t) = 0$ keine gleichen Wurzeln hat, daß zwar $F(r,\varrho)$, aber keiner von den anderen Faktoren $F(r,\varrho^m)$ verschwindet. Die Gleichungen

\weq{18}{F(r,t) = 0, \quad \varPhi(t) = 0}

haben daher nur die eine Wurzel $t = \varrho$ miteinander gemein, und durch Aufsuchen ihres größten gemeinschaftlichen Teilers \gesperrt{findet man $\varrho$ rational ausgedrückt durch $r$, d.\,h. durch die Wurzeln der Teilungsgleichung.}

Diese Ausdrücke für $\varrho$ ändern ihren Wert nicht, wenn auf $r$ eine Substitution der Gruppe $\mathfrak{B}$ angewandt wird, während [nach (12)] $\varrho$ durch eine Substitution der Gruppe $\mathfrak{A}$ in eine Potenz von $\varrho$ übergeht.

Die \gesperrt{Abelschen} Relationen zeigen, daß durch die in $\mathfrak{A}$ enthaltene Substitution

\weq{19}{\begin{pmatrix} m, & 0 \\ 0, & 1 \end{pmatrix}}

$\varrho$ in $\varrho^m$ übergeht; denn setzen wir

\[
e^{\frac{8\pi i}{n}} = \varrho^\lambda,
\]

so lautet eine der \gesperrt{Abelschen} Relationen

\weq{20}{\overset{\nu}{\Sigma} \varrho^{\lambda\nu\nu'} x_{\nu,\nu'} = 0,}

worauf, wenn für $\varrho$ der Ausdruck durch $r$ gesetzt wird, \gesperrt{alle} Substitutionen von $\mathfrak{A}$, also auch (19), anwendbar sind. Nach § 62, (12) bleibt aber (20) bei dieser Substitution \gesperrt{nur} richtig, wenn $\varrho$ in $\varrho^m$ übergeht.

Da $\varrho$ durch die Substitutionen in $\mathfrak{B}$ nicht geändert wird, und da die Gruppe $\mathfrak{A}$ aus $\mathfrak{B}$ durch Zusammensetzung mit (19) entsteht, so folgt, daß durch irgend eine Substitution in $\mathfrak{A}$

\[
\begin{pmatrix} a, & b \\ c, & \partial \end{pmatrix}
\]
%CONT>>
